% ---------------------------------------------------
% Trabajo Final de Grado
% Author: Fabián González Lence <alu0101549491@ull.edu.es>
% Chapter: Modelos de IA (versión reducida)
% ----------------------------------------------------

\cleardoublepage
\chapter{Modelos de IA} \label{chap:AIModels}

Este capítulo describe los modelos y herramientas de \ac{IA} empleados a lo largo del trabajo, justifica su selección y documenta el cambio de enfoque metodológico introducido a partir del cuarto proyecto, como se explica en la Sección~\ref{sec:approach-change}. Cada herramienta fue asignada a un rol concreto dentro del flujo de desarrollo, configurando un equipo de agentes especializados que no fue fijo desde el inicio, sino que evolucionó con la experiencia acumulada en cada proyecto hasta encontrar la combinación más adecuada para cada fase y nivel de complejidad.

\section{Panorama actual de los LLMs en el desarrollo de software} \label{sec:llm-overview}

En el momento de realización de este trabajo, el ecosistema de \ac{LLMs} orientados al desarrollo de software ha alcanzado un grado de madurez suficiente para abordar tareas de complejidad real, más allá de la generación de fragmentos de código aislados.
Modelos de propósito general como GPT-5 (OpenAI)\footnote{ \href{https://cdn.openai.com/gpt-5-system-card.pdf}{GPT-5 System Card -- \texttt{cdn.openai.com/gpt-5-system-card.pdf}}}, Claude Sonnet 4.5 (Anthropic)\footnote{\href{https://www.anthropic.com/claude-sonnet-4-5-system-card}{Claude Sonnet 4.5 System Card -- \texttt{anthropic.com/claude-sonnet-4-5-system-card}}}, Gemini 3 Pro (Google)\footnote{\href{https://storage.googleapis.com/deepmind-media/Model-Cards/Gemini-3-Pro-Model-Card.pdf}{Gemini 3 Model Card \\ \hspace*{1.9em} \texttt{storage.googleapis.com/deepmind-media/Model-Cards/Gemini-3-Pro-Model-Card.pdf}}} o Mistral Large (Mistral AI)\footnote{\href{https://docs.mistral.ai/models/model-cards/mistral-large-2-0-24-07}{Mistral Large Model Card \\ \hspace*{1.9em} \texttt{docs.mistral.ai/models/model-cards/mistral-large-2-0-24-07}}} compiten en tareas de razonamiento sobre código, diseño arquitectónico y generación de documentación técnica.
Paralelamente, herramientas en el \ac{IDE} como GitHub Copilot han evolucionado desde el autocompletado inteligente hasta un modo agente capaz de ejecutar flujos de trabajo completos de forma autónoma dentro del mismo.

Esta dualidad entre modelos conversacionales y agentes integrados resultó especialmente relevante para el presente trabajo: los primeros ofrecen flexibilidad para el diseño y la supervisión de alto nivel, mientras que los agentes operan directamente sobre el sistema de ficheros, ejecutan comandos, leen errores del compilador y realizan múltiples iteraciones de corrección sin salir del editor. Conviene señalar que el panorama de los \ac{LLMs} evoluciona muy rápidamente; por ello, el análisis de este capítulo no pretende ofrecer una clasificación definitiva de modelos, sino documentar con fidelidad las herramientas concretas utilizadas, el rol asignado a cada una y el comportamiento observado en este experimento.

\section{Modelos y herramientas utilizados} \label{sec:models-description}

La configuración del equipo de \ac{IA} no fue el resultado de una elección única al inicio del trabajo, sino de un proceso iterativo en el que la experiencia acumulada en cada proyecto ajustó qué herramienta resultaba más adecuada para cada fase y nivel de complejidad.
Cada modelo fue asignado al rol en que su fortaleza diferencial fuera más aprovechable y su debilidad tuviera menor impacto.

\subsection{Claude} \label{subsec:claude}

Claude\footnote{\href{https://claude.ai}{Claude (\textit{chatbot}) -- \texttt{claude.ai}}} fue el designado como \ac{IA} Arquitecto en los cinco proyectos. Se utilizó principalmente el modelo Claude Sonnet 4.5, disponible a través del propio chatbot \textit{claude.ai} y como agente integrado en GitHub Copilot. Su elección para el rol arquitectónico se fundamentó en tres fortalezas: una ventana de contexto amplia que le permitía procesar diagramas de clases y especificaciones extensas en una única interacción; una notable capacidad para generar documentación técnica detallada a partir de descripciones en lenguaje natural; y una coherencia arquitectónica que se mantiene consistente a lo largo de conversaciones largas.

La tarea concreta de Claude en cada proyecto consistía en, partiendo de las especificaciones de requisitos y los diagramas de clases y casos de uso, generar la estructura completa del proyecto: árbol de directorios, ficheros de configuración, esqueletos de clases con métodos declarados sin implementar y documentación base. A partir del cuarto proyecto, Claude continuó en este rol integrado como agente personalizado dentro de GitHub Copilot; adicionalmente, durante el proyecto \ac{TENNIS} se empleó Claude Opus 4.6 en fases de mayor exigencia de razonamiento.

\subsection{Mistral AI} \label{subsec:mistral}
Mistral AI\footnote{\href{https://chat.mistral.ai/chat}{Mistral (\textit{chatbot}) -- \texttt{chat.mistral.ai/chat}}} desempeñó el rol de \ac{IA} Desarrolladora durante los tres primeros proyectos (\ac{HANGMAN}, \ac{PLAYER} y \ac{BALATRO}), siendo responsable de la implementación del código fuente de cada módulo a partir de los artefactos de diseño obtenidos de Claude. Su elección respondía a su buena relación entre capacidad de generar código, accesibilidad y la novedad que supuso cuando salió al mercado.

En la práctica, Mistral produjo implementaciones funcionales para módulos de complejidad media; sin embargo, a medida que los proyectos escalaron, se hicieron evidentes sus limitaciones: la gestión del estado entre módulos interdependientes resultaba inconsistente al superar ciertos límites de contexto, y la generación de patrones arquitectónicos avanzados requería numerosas iteraciones y correcciones. Estas limitaciones motivaron el cambio de enfoque descrito en la Sección~\ref{sec:approach-change}.

\subsection{GitHub Copilot} \label{subsec:copilot}

GitHub Copilot\footnote{\href{https://github.com/copilot/agents}{GitHub Copilot (\textit{chatbot}) -- \texttt{github.com/copilot/agents}}} asumió desde el inicio el rol de agente \ac{IA} Revisora en todos los proyectos y, en los dos últimos proyectos, (\ac{CARTO} y \ac{TENNIS}), también los roles de Desarrolladora y \textit{Tester}, consolidando en una única herramienta las responsabilidades que anteriormente se distribuían entre varias.

Su principal fortaleza frente al enfoque conversacional previo reside en la integración directa con el sistema de ficheros, gracias a su extensión de Visual Studio Code\footnote{\href{https://code.visualstudio.com/docs/copilot/overview}{GitHub Copilot en Visual Studio Code -- \texttt{code.visualstudio.com/docs/copilot/overview}}}: el agente lee el código existente, ejecuta comandos de compilación y estilizado (\textit{linting}), interpreta los errores y aplica correcciones iterativas sin salir del editor. Esta capacidad, sumada a la posibilidad de definir agentes personalizados mediante ficheros de configuración en el repositorio, permitió mapear cada agente a uno de los cuatro roles ---Arquitectura, Codificación, Revisión y \textit{Testing}--- con reglas y modelos específicos para cada fase.

En cuanto al modelo subyacente, se empleó en una primera instancia Claude Sonnet 4.5 para los roles de Arquitecto y Desarrolladora, y modelos de la familia GPT-5 (mini, 5.4) para los roles de Revisora y \textit{Tester}, utilizando de forma puntual modelos más potentes como Claude Opus 4.6.

\subsection{Qwen AI} \label{subsec:qwen}

Qwen AI\footnote{\href{https://qwen.ai/home}{Qwen AI (\textit{chatbot}) -- \texttt{qwen.ai/home}}} fue el agente de \ac{IA} \textit{Tester} durante los tres primeros proyectos, responsable de generar los conjuntos de pruebas unitarias con Jest. 
Su elección respondía a su eficacia en la generación de casos de prueba siguiendo el patrón \ac{AAA} \cite{Wei:2024:AAA}, con una tasa de compilación inicial elevada y cobertura de casos normales, límite y excepcionales.
Al igual que Mistral, fue desplazado a partir de \ac{CARTO} por la inviabilidad del flujo manual de copiar código entre interfaces chatbot para proyectos de mayor escala.

\subsection{ChatGPT} \label{subsec:chatgpt}

ChatGPT\footnote{\href{https://chatgpt.com}{ChatGPT (\textit{chatbot}) -- \texttt{chatgpt.com}}} fue utilizado como herramienta de apoyo en las fases iniciales de los primeros proyectos, principalmente para la redacción de especificaciones de requisitos y la exploración de alternativas arquitectónicas.
Su fortaleza resultó ser su fluidez en la generación de texto técnico; sin embargo, en tareas de codificación extensa mostró tendencia a perder coherencia contextual en conversaciones largas, lo que limitó su utilidad.
Algunos modelos GPT pudieron ser aprovechados en las fases de revisión y pruebas  a través de GitHub Copilot gracias a que pudieron utilizarse de forma integrada en el \ac{IDE}.

\subsection{Google Gemini y Perplexity --- Uso puntual} \label{subsec:others}

Google Gemini\footnote{\href{https://gemini.google.com/app}{Google Gemini (\textit{chatbot}) -- \texttt{gemini.google.com}}} fue utilizado de forma puntual en el proyecto \ac{BALATRO} para la generación de los recursos (\textit{assets}) gráficos correspondientes a algunas de las cartas del juego, resolviendo directamente una necesidad de producción visual sin requerir diseño manual ni intervención humana. Perplexity\footnote{\href{https://www.perplexity.ai}{Perplexity (\textit{chatbot}) -- \texttt{perplexity.ai}}} fue evaluado en las fases iniciales como motor de búsqueda potenciado por \ac{IA}, pero no llegó a integrarse en el flujo de trabajo al no ofrecer ventajas respecto a las herramientas adoptadas.

\section{El cambio de enfoque metodológico} \label{sec:approach-change}

La transición de los proyectos 1--3 al proyecto 4 (\ac{CARTO}) supuso un punto de inflexión, motivado por las limitaciones de escalabilidad del flujo conversacional hasta entonces empleado. En los tres primeros proyectos, la \ac{IA} Desarrolladora operaba como \textit{chatbot} web: se transfería el código al editor módulo a módulo de forma manual. Este flujo resultaba viable para proyectos de 9 a 15 módulos, pero mostró tres limitaciones críticas al escalar: pérdida de coherencia inter-módulo al generar ficheros en conversaciones independientes; fricción operativa del ciclo manual de copiar, pegar, detectar errores y volver al \textit{chatbot}; e imposibilidad práctica de revisar y probar a la escala de proyectos como \ac{CARTO}, que llegó a gestionar aproximadamente 240 ficheros en cuatro capas de arquitectura. La solución fue adoptar GitHub Copilot en modo agente como herramienta principal para Codificación, Revisión y \textit{Testing} a partir de \ac{CARTO}, decisión influenciada por un curso de agentes personalizados de GitHub Copilot cursado en el ámbito de ``\textit{Prácticas Externas}''.

Los \textit{prompts} de los tres roles fueron rediseñados para el modelo agente. La Revisora adoptó una metodología de tres etapas: exploración de la base de código y generación de una lista de tareas; revisión sistemática por grupos de ficheros con informe de incidencias clasificadas por severidad; e implementación de correcciones siguiendo el informe. El \textit{Tester} generó pruebas \textit{end-to-end} con Playwright a partir de los diagramas de casos de uso: en una primera pasada se definían los escenarios y en una segunda se implementaban. En \ac{TENNIS}, la Desarrolladora incorporó además una fase de planificación previa que categoriza los módulos y establece su orden de dependencia antes de comenzar la implementación.

\section{Comparativa de modelos} \label{sec:comparativa-modelos}

La Tabla~\ref{tab:comparativa-modelos} sintetiza los modelos empleados: en P1--P3 el equipo se distribuía entre tres \textit{chatbots} independientes; en P4--P5 los roles convergieron en GitHub Copilot con modelos de Claude y GPT como subyacentes. Ningún modelo resultó adecuado para todas las fases, lo que justifica el enfoque multiagente; el rol de Revisora fue el único constante, evidenciando que la integración con el sistema de ficheros resulta crítica en las tareas de revisión.

\begin{table}[H]
\centering
\small
\renewcommand{\arraystretch}{1.2}
\begin{tabularx}{\textwidth}{|
>{\raggedright\arraybackslash}p{2cm}|
>{\raggedright\arraybackslash}p{3cm}|
>{\raggedright\arraybackslash}p{3.9cm}|
Y|
}
\hline
\textbf{Modelo} & \textbf{Rol} & \textbf{Fortalezas} & \textbf{Limitaciones} \\
\hline
\textbf{Claude} &
Arquitecto (P1--P4) &
Contexto amplio; gran coherencia; fidelidad de diseño &
Sin integración directa; gran coste en modelos de mayor capacidad \\
\hline
\textbf{Mistral} &
Desarrolladora (P1--P3) &
Buen rendimiento en módulos de complejidad media &
Pérdida de coherencia inter-módulo; dificultad con patrones avanzados \\
\hline
\textbf{GitHub Copilot} &
Revisora (P1--P5); Desarrolladora y \textit{Tester} (P4--P5); Arquitecto (P5) &
Integración directa; ejecución iterativa de comandos; gestión coherente a gran escala &
Uso de modelos externos con consumo de \textit{tokens} moderado \\
\hline
\textbf{Qwen} &
\textit{Tester} (P1--P3) &
Eficaz con patrón \ac{AAA}; cobertura de un amplio rango de casos &
Flujo manual no escalable en grandes proyectos \\
\hline
\textbf{ChatGPT} &
Apoyo en requisitos (P1--P3) &
Fluidez técnica; familiaridad con convenciones estándar &
Pérdida de coherencia en conversaciones largas \\
\hline
\textbf{Perplexity} &
Consulta puntual (ninguno) &
Búsqueda en tiempo real con referencias &
No aporta ventajas en desarrollo de código \\
\hline
\end{tabularx}
\caption{Comparativa de modelos y herramientas de \ac{IA} empleados. P1-\ac{HANGMAN}, P2-\ac{PLAYER}, P3-\ac{BALATRO}, P4-\ac{CARTO} y P5-\ac{TENNIS}.}
\label{tab:comparativa-modelos}
\end{table}

\section{Uso de IA en la redacción de este documento} \label{sec:ia-memoria}

La redacción de esta memoria se apoyó en el modelo Claude Sonnet 4.6 como asistente de escritura LaTeX\footnote{\href{https://www.latex-project.org}{LaTeX -- \texttt{latex-project.org}}}, con un uso de asistencia y transformación, no generativo desde cero: el punto de partida se establece en notas propias ---borradores en lenguaje llano, listas de ideas y registros en un cuaderno Notion\footnote{\href{https://www.notion.com/es-es}{Notion -- \texttt{notion.com/es-es}}} personal--- y el agente proponía reescrituras adaptadas al registro académico del \ac{TFG}. Las tareas delegadas abarcaron la estructuración de capítulos, la condensación de ideas para eliminar redundancias y la aplicación de convenciones de formato LaTeX ---italización de términos en inglés, uso consistente de acrónimos, ajuste de tablas---. En ningún caso operó como fuente primaria de información.